Synthetic face imagery produced by generative adversarial networks, face-swapping tools and diffusion models can now be hard to
distinguish from genuine photographs, which creates a demand for forensic tools that give not only a verdict but also inspectable evidence.
This report presents \textbf{Verif-AI}, an end-to-end framework for face-based deepfake detection and forensic analysis that works on both
images and videos. Faces are located with MTCNN and passed to a two-stream convolutional classifier: an EfficientNet-B0 stream that
reads the RGB face, and a second EfficientNet-B0 stream that reads a noise residual produced by a fixed three-kernel Spatial Rich Model
(SRM) filter bank. The two feature vectors are fused by a small classification head into a single raw sigmoid score. Gradient-weighted
Class Activation Mapping (Grad-CAM) shows which facial regions influenced the RGB stream's score, a video pipeline adds face tracking,
frame-level timelines, suspicious time ranges and stored evidence frames, and a PDF forensic report records the result together with the
model checkpoint hash and the limitations. The system is implemented as a FastAPI backend and a React/TypeScript dashboard in which
uploads are deleted immediately after analysis and stored results can be deleted by the user.

The evaluation is reported as it was obtained, including its negative findings. The selected production model (v2, trained on a sample of
the DeepFakeFace dataset) reaches an accuracy of 0.875 and a ROC-AUC of 0.941 on held-out identities seen-generator data, but only 0.746 and 0.849 on an
unseen text-to-image generator, and it falsely flags genuine video from the Celeb-DF dataset at chance level. Diagnostic experiments
traced this failure to an image-quality shortcut in the training data, and adding Celeb-DF video frames and quality augmentation (Model 4)
raised the video-level accuracy on held-out Celeb-DF identities to 0.980, which is within-dataset evidence only. A further experiment
with modern-generator faces (E3) detected held-out fakes of those generators almost perfectly, but without a generator-disjoint test it
cannot be read as evidence of general detection, and it did not meet the pre-registered promotion rule. Scores are raw sigmoid outputs,
not calibrated probabilities. The latest model, E4, is the default of the local examiner demonstration; it did not pass the promotion rule, so the public deployment still serves v2. The report does not claim a general detector of AI-generated images or videos, and the public deployment on
free hosting tiers runs in a restricted mode that is not production-ready.

\medskip\noindent\textbf{Keywords:} deepfake detection, two-stream network, EfficientNet, SRM noise residuals, Grad-CAM, digital forensics, FastAPI, React.
